\documentclass[11pt,a4paper]{article}

\usepackage{fontspec}
\IfFontExistsTF{Times New Roman}{\setmainfont{Times New Roman}}{\setmainfont{Tinos}}
\IfFontExistsTF{Consolas}{\setmonofont{Consolas}[Scale=0.92]}{\setmonofont{Liberation Mono}[Scale=0.92]}
\usepackage[english]{babel}
\usepackage{geometry}
\usepackage{xcolor}
\usepackage{enumitem}
\usepackage{listings}
\usepackage{booktabs}
\usepackage{array}
\usepackage{amsmath,amssymb}
\usepackage{titlesec}
\usepackage{fancyhdr}
\usepackage{parskip}
\usepackage{tcolorbox}
\tcbuselibrary{skins}
\usepackage{tabularx}
\usepackage{lastpage}
\usepackage{hyperref}

\geometry{top=1.5cm,bottom=1.5cm,left=1.7cm,right=1.7cm}

\definecolor{primaryblue}{RGB}{18,65,115}
\definecolor{secondaryteal}{RGB}{14,102,110}
\definecolor{SoftBlue}{RGB}{232,242,250}
\definecolor{CodeBg}{RGB}{248,250,252}
\definecolor{DarkGray}{RGB}{55,55,55}
\definecolor{boxbg}{RGB}{248,250,252}
\definecolor{borderblue}{RGB}{150,185,215}

\setlength{\parindent}{0pt}
\setlength{\parskip}{0pt}
\setlist[itemize]{leftmargin=1.3em,itemsep=1pt,topsep=1pt,parsep=0pt,partopsep=0pt}
\setlist[enumerate]{leftmargin=1.5em,itemsep=1.5pt,topsep=1pt,parsep=0pt,partopsep=0pt}
\setlist[enumerate,2]{topsep=0.5pt,itemsep=0pt,parsep=0pt,partopsep=0pt}

\lstset{
  language=Python,
  basicstyle=\ttfamily\scriptsize,
  keywordstyle=\color{blue!75!black}\bfseries,
  stringstyle=\color{red!65!black},
  commentstyle=\color{primaryblue!80!black}\itshape,
  showstringspaces=false,
  columns=fullflexible,
  keepspaces=true,
  breaklines=true,
  frame=single,
  rulecolor=\color{primaryblue!45},
  backgroundcolor=\color{CodeBg},
  xleftmargin=0.2em,
  xrightmargin=0.2em,
  aboveskip=0.2em,
  belowskip=0.2em,
  tabsize=4
}

\pagestyle{fancy}
\fancyhf{}
\setlength{\headheight}{14pt}
\fancyhead[L]{\footnotesize\textbf{DSAI1003 -- Fundamental Programming Concepts in Python}}
\fancyhead[R]{\footnotesize\textbf{Decision Making Exam -- 60 Minutes}}
\fancyfoot[L]{\scriptsize Faculty of Data Science and Artificial Intelligence -- National Economics University}
\fancyfoot[R]{\footnotesize Page \thepage\ of 6}
\renewcommand{\headrulewidth}{0.4pt}
\renewcommand{\footrulewidth}{0.3pt}

\titleformat{\section}{\large\bfseries\color{primaryblue}}{}{0pt}{}
\titleformat{\subsection}{\normalsize\bfseries\color{primaryblue}}{}{0pt}{}
\titlespacing*{\section}{0pt}{3.5pt}{1.5pt}
\titlespacing*{\subsection}{0pt}{2.5pt}{1pt}

\newtcolorbox{headerbox}{
  colback=white,
  colframe=primaryblue,
  boxrule=0.7pt,
  arc=1.5mm,
  left=2.5mm,
  right=2.5mm,
  top=1.2mm,
  bottom=1.2mm
}

\newtcolorbox{answerbox}[2][]{
  enhanced,
  colback=white,
  colframe=DarkGray!45,
  boxrule=0.5pt,
  arc=1mm,
  left=2mm,
  right=2mm,
  top=1.2mm,
  bottom=1.2mm,
  title=\footnotesize\color{DarkGray}\textbf{#2},
  #1
}

\hypersetup{
  colorlinks=true,
  linkcolor=primaryblue,
  urlcolor=primaryblue,
  pdftitle={DSAI1003 - Decision Making in Python Exam},
  pdfauthor={Faculty of Data Science and Artificial Intelligence - National Economics University (NEU)}
}

\begin{document}

% =========================================================================
% PAGE 1: OFFICIAL HEADER + ANSWER SHEET + SECTION A (Q1 - Q5)
% =========================================================================

\noindent
\begin{minipage}[t]{0.64\textwidth}
  \raggedright\small
  \textbf{NATIONAL ECONOMICS UNIVERSITY}\\
  \textbf{COLLEGE OF TECHNOLOGY}\\
  \textbf{FACULTY OF DATA SCIENCE \& ARTIFICIAL INTELLIGENCE (FDA)}
\end{minipage}%
\hfill
\begin{minipage}[t]{0.34\textwidth}
  \raggedleft\footnotesize
  \textbf{OFFICIAL MIDTERM EXAMINATION}\\
  Academic Year: 2026--2027\\
  Semester: Fall 2026 \quad Course Code: \textbf{DSAI1003}
\end{minipage}

\vspace{0.15em}

\begin{center}
  {\Large\bfseries\color{primaryblue} FUNDAMENTAL PROGRAMMING CONCEPTS IN PYTHON}\\[0.08em]
  {\normalsize\bfseries PROGRESS EXAM: DECISION MAKING \& BOOLEAN LOGIC}\\[0.08em]
  {\normalsize\textbf{Paper Code:} \rule{1.6cm}{0.4pt} \quad | \quad \textbf{Duration: 60 Minutes} (Excluding paper distribution)}
\end{center}

\vspace{0.05em}

\begin{headerbox}
\begin{tabular*}{\textwidth}{@{\extracolsep{\fill}}ll@{}}
  \textbf{Candidate Full Name:} \dotfill & \textbf{Student ID (MSSV):} \dotfill \\[0.2em]
  \textbf{Class / Cohort:} \dotfill & \textbf{Exam Room / Seat No.:} \dotfill
\end{tabular*}

\vspace{0.2em}
\renewcommand{\arraystretch}{1.05}
\setlength{\tabcolsep}{4.5pt}
\centering
\begin{tabular}{|l|c|c|c|c||c|c|}
\hline
\textbf{Section} & \textbf{A} & \textbf{B} & \textbf{C} & \textbf{D} & \textbf{Total (100)} & \textbf{Examiner Signature} \\
\hline
\textbf{Max Score} & 20 & 25 & 25 & 30 & \textbf{100} & \\[0.3em]
\cline{1-6}
\textbf{Score} & & & & & & \\[0.4em]
\hline
\end{tabular}

\vspace{0.15em}
\raggedright
\scriptsize
\textbf{Instructions:} Closed-book examination. All electronic devices and internet access are prohibited. Answer all questions directly in the designated spaces. In programming tasks, apply conditional statements (\texttt{if}, \texttt{elif}, \texttt{else}), relational/Boolean operators, and input validation. Do \textbf{not} use loops or custom functions (\texttt{def}) unless explicitly requested. Assume users enter data in expected types unless validation is specifically instructed. Do not detach pages.
\end{headerbox}

\vspace{1mm}
\noindent
\textbf{MULTIPLE-CHOICE ANSWER SHEET} \textit{(Write A, B, C, or D in the corresponding box)}:
\begin{center}
\setlength{\tabcolsep}{5pt}
\renewcommand{\arraystretch}{1.0}
\begin{tabular}{|c|c|c|c|c|c|c|c|c|c|c|}
\hline
\textbf{Question} & \textbf{1} & \textbf{2} & \textbf{3} & \textbf{4} & \textbf{5} & \textbf{6} & \textbf{7} & \textbf{8} & \textbf{9} & \textbf{10} \\
\hline
\textbf{Answer} & & & & & & & & & & \\[1.8mm]
\hline
\end{tabular}
\end{center}
\vspace{-2.5mm}
\section*{Section A: Decision Making \& Boolean Logic Concepts \hfill \normalsize\normalfont(20 points)}
\textit{Answer all 10 questions. Each correct answer carries 2 points. Suggested time: 12 minutes.}

\begin{enumerate}[label=\textbf{Q\arabic*.},leftmargin=2em,itemsep=1.5pt,topsep=1pt]

\item Which statement correctly describes compound statements and code blocks in Python?
\begin{enumerate}[label=\Alph*.,leftmargin=2em]
  \item A colon (\texttt{:}) is optional after the \texttt{if} condition as long as the block is indented.
  \item Every branch block must be consistently indented; mixing tabs and spaces raises an \texttt{IndentationError}.
  \item Python uses curly braces \texttt{\{\}} to delimit blocks; indentation is solely for readability.
  \item An \texttt{else} clause can exist independently without a preceding \texttt{if} or \texttt{elif} statement.
\end{enumerate}

\item What happens when executing the statement \texttt{if x = 10:} in Python?
\begin{enumerate}[label=\Alph*.,leftmargin=2em]
  \item It assigns \texttt{10} to \texttt{x} and executes the branch because \texttt{10} is considered truthy.
  \item It checks if \texttt{x} equals \texttt{10} and returns a boolean value.
  \item It causes a \texttt{SyntaxError} because \texttt{=} is assignment, not the equality operator \texttt{==}.
  \item It automatically converts \texttt{x} to an integer before comparing.
\end{enumerate}

\item Why is the direct comparison \texttt{if x == 0.3:} hazardous when \texttt{x = 0.1 + 0.2}, and what is the standard fix?
\begin{enumerate}[label=\Alph*.,leftmargin=2em]
  \item Relational operators cannot compare floats; one must convert to strings using \texttt{str(x) == "0.3"}.
  \item Binary representation causes roundoff error; one should verify \texttt{abs(x - 0.3) < 1E-14}.
  \item Floating-point values are automatically rounded to integers inside conditional statements.
  \item One must always use the identity operator instead: \texttt{if x is 0.3:}.
\end{enumerate}

\item In Python, what is the boolean value of the expression \texttt{"apple" < "Banana"} and what is the reason?
\begin{enumerate}[label=\Alph*.,leftmargin=2em]
  \item \texttt{True}, because \texttt{"apple"} precedes \texttt{"Banana"} in standard alphabetical order.
  \item \texttt{False}, because uppercase letters have smaller Unicode values than lowercase letters (\texttt{'B'} is 66, \texttt{'a'} is 97).
  \item \texttt{True}, because \texttt{"apple"} has fewer characters than \texttt{"Banana"}.
  \item A \texttt{TypeError} is raised because strings cannot be compared with relational operators.
\end{enumerate}

\item Given \texttt{score = 85} in \texttt{if score >= 60: g = "D"} followed by \texttt{elif score >= 80: g = "B"} and \texttt{elif score >= 90: g = "A"}, what is \texttt{g}?
\begin{enumerate}[label=\Alph*.,leftmargin=2em]
  \item \texttt{"B"}, because $85 \ge 80$ is the most specific matching range.
  \item \texttt{"D"}, because \texttt{score >= 60} is evaluated first and evaluates to \texttt{True}; remaining branches are skipped.
  \item \texttt{"A"}, because Python evaluates \texttt{elif} branches in descending numerical order.
  \item \texttt{"F"}, because matching multiple conditions produces an ambiguous branch error.
\end{enumerate}

\end{enumerate}

\newpage

% ==============================================================================
% PAGE 2: SECTION A (Q6 - Q10) + SECTION B Q11
% ==============================================================================

\begin{enumerate}[label=\textbf{Q\arabic*.},leftmargin=2em,itemsep=1.5pt,topsep=1pt,resume]

\item Following Python's operator precedence rules, what is the truth value of \texttt{not False and False or True}?
\begin{enumerate}[label=\Alph*.,leftmargin=2em]
  \item \texttt{True}, because \texttt{not} evaluates first (\texttt{True}), then \texttt{and} (\texttt{False}), then \texttt{or} (\texttt{True}).
  \item \texttt{False}, because \texttt{and} has higher precedence than \texttt{not}.
  \item \texttt{False}, because boolean operators are evaluated strictly from left to right.
  \item A \texttt{TypeError} occurs because parentheses are required when mixing logical operators.
\end{enumerate}

\item Given \texttt{x = 0}, what is the outcome of evaluating the expression \texttt{x != 0 and (10 / x > 2)}?
\begin{enumerate}[label=\Alph*.,leftmargin=2em]
  \item A \texttt{ZeroDivisionError} is raised because \texttt{10 / x} is evaluated.
  \item \texttt{False}, because the first operand is \texttt{False}, so Python short-circuits and skips the division.
  \item \texttt{True}, because Python treats division by zero as positive infinity.
  \item \texttt{None}, because the expression terminates abnormally without a return value.
\end{enumerate}

\item According to De Morgan's Laws, which condition is logically equivalent to \texttt{not (status == "VIP" or points >= 100)}?
\begin{enumerate}[label=\Alph*.,leftmargin=2em]
  \item \texttt{status != "VIP" or points < 100}
  \item \texttt{status != "VIP" and points < 100}
  \item \texttt{not (status == "VIP") or not (points >= 100)}
  \item \texttt{status == "VIP" and points < 100}
\end{enumerate}

\item Which string method checks whether all characters in string \texttt{s} are digits and \texttt{s} has at least one character?
\begin{enumerate}[label=\Alph*.,leftmargin=2em]
  \item \texttt{s.isnumeric\_int()}
  \item \texttt{s.count()}
  \item \texttt{s.isdigit()}
  \item \texttt{s.all\_digits()}
\end{enumerate}

\item What is the minimum number of test cases required to achieve full branch coverage for an \texttt{if-elif-elif-else} construct?
\begin{enumerate}[label=\Alph*.,leftmargin=2em]
  \item 1 test case
  \item 3 test cases
  \item 4 test cases (one for \texttt{if}, one for each of the two \texttt{elif} blocks, and one for \texttt{else})
  \item 8 test cases
\end{enumerate}

\end{enumerate}

\section*{Section B: Output Tracing, Logic Diagnosis \& Error Correction \hfill \normalsize\normalfont(25 points)}
\textit{Answer all 5 questions. Each question carries 5 points. Suggested time: 15 minutes.}

\subsection*{Q11. Exact Output Tracing with Nested Decisions \& Surcharges (5 points)}
Predict the \textbf{exact screen output}, including punctuation, spaces, and line breaks:

\begin{lstlisting}
age = 19
is_student = True
base_fare = 40.0
discount = 0.0

if age < 12:
    discount = 0.50
elif age < 25:
    if is_student:
        discount = 0.30
    else:
        discount = 0.15
elif age >= 65:
    discount = 0.40
else:
    discount = 0.0

final_fare = base_fare * (1.0 - discount)
print(f"Base: ${base_fare:.2f} | Discount: {discount * 100:.0f}%")
print(f"Payable: ${final_fare:.2f}")
print(f"Eligible: {discount >= 0.30 and age >= 18}")
\end{lstlisting}

\begin{answerbox}[height=3.0cm]{Output Box for Q11}
\end{answerbox}

\newpage

% ==============================================================================
% PAGE 3: SECTION B, Q12--Q15
% ==============================================================================

\subsection*{Q12. Evaluation of Boolean Expressions and Short-Circuiting (5 points)}
Write the exact boolean result (\texttt{True} or \texttt{False}) for each expression:

\begin{enumerate}[label=\alph*),leftmargin=2em,itemsep=2pt]
  \item \texttt{10 < 20 <= 25}: \underline{\hspace{4cm}}
  \item \texttt{"py" in "Python".lower() and not (5 > 8)}: \underline{\hspace{4cm}}
  \item \texttt{False or not False and not True}: \underline{\hspace{4cm}}
  \item \texttt{"cat" > "Cat" and "10" < "2"}: \underline{\hspace{4cm}}
  \item \texttt{len("neu") == 3 or (1 / 0 == 1)}: \underline{\hspace{4cm}}
\end{enumerate}

\subsection*{Q13. Diagnose and Correct Common Decision Logic Errors (5 points)}
For each code scenario, state the underlying defect and provide the corrected Python statement.

\renewcommand{\arraystretch}{1.3}
\begin{tabularx}{\textwidth}{|p{5.1cm}|X|X|}
\hline
\textbf{Code / Situation} & \textbf{Problem / Logic Flaw} & \textbf{Correction / Valid Code} \\
\hline
\texttt{choice = input("Option: ")}\\
\texttt{if choice == "A" or "B":}\\
\texttt{\ \ \ \ print("Accepted")} & & \\[3.5mm]
\hline
\texttt{x = 0.1 + 0.2}\\
\texttt{if x == 0.3:}\\
\texttt{\ \ \ \ print("Equal")} & & \\[3.5mm]
\hline
\texttt{val = input("Enter int: ")}\\
\texttt{if int(val) > 0:}\\
(Risk: user enters \texttt{"abc"}) & & \\[3.5mm]
\hline
\end{tabularx}

\subsection*{Q14. String Analysis Methods and Substring Inspection (5 points)}
Predict the exact screen output of the following verification routine:

\begin{lstlisting}
code = "NEU-DSFE2026"
is_prefix = code.startswith("NEU")
has_dash = "-" in code
suffix_digits = code[-4:].isdigit()
is_clean = code.isalnum()
print(f"Prefix: {is_prefix}, Dash: {has_dash}")
print(f"Suffix Digits: {suffix_digits}")
print(f"Alphanumeric: {is_clean}")
\end{lstlisting}

\begin{answerbox}[height=2.35cm]{Output Box for Q14}
\end{answerbox}

\subsection*{Q15. De Morgan's Laws and Boundary Test Cases (5 points)}
\begin{itemize}[leftmargin=2em]
  \item Apply De Morgan's Law to rewrite the condition without the outer \texttt{not}:\\
  Original: \texttt{not (temp >= 18 and temp <= 26)}\\
  Rewritten: \underline{\hspace{10cm}}
  \item State the \textbf{critical boundary values} that must be tested to thoroughly verify \texttt{18 <= temp <= 26}: \hrulefill
\end{itemize}

\newpage

% ==============================================================================
% PAGE 4: SECTION C
% ==============================================================================

\section*{Section C: Program Design \& Logic Specification \hfill \normalsize\normalfont(25 points)}
\textit{Answer all tasks. Suggested time: 15 minutes.}

\begin{tcolorbox}[colback=boxbg,colframe=secondaryteal,arc=1.5mm,boxrule=0.6pt,top=1.5mm,bottom=1.5mm]
\textbf{Problem Specification: Air Cargo \& Express Courier Pricing Engine}\\[0.5mm]
\small
A logistics platform computes parcel delivery cost based on weight $w$ (kg), service speed (\texttt{"S"} for Standard, \texttt{"E"} for Express), and corporate membership status (\texttt{is\_member}):
\begin{itemize}[itemsep=0.5pt,topsep=0.5pt]
  \item \textbf{Input Validation:} Weight must satisfy $0 < w \le 50.0$; service must be \texttt{"S"} or \texttt{"E"}. Otherwise, reject as invalid.
  \item \textbf{Base Tiers:} Tier 1 ($0 < w \le 2$ kg): flat \$10.00. Tier 2 ($2 < w \le 10$ kg): \$10.00 + \$3.50/kg for weight over 2.0 kg. Tier 3 ($10 < w \le 50$ kg): \$38.00 + \$2.00/kg for weight over 10.0 kg.
  \item \textbf{Express Surcharge:} If \texttt{service == "E"}, add 40\% surcharge on base cost (minimum surcharge \$8.00).
  \item \textbf{Member Discount:} If \texttt{is\_member} is \texttt{True} and subtotal exceeds \$30.00, deduct a \$5.00 loyalty credit.
\end{itemize}
\end{tcolorbox}

\vspace{0.5mm}
\textbf{Task 1: Identify Decision Criteria, Inputs \& Outputs} (6 points)
\begin{itemize}[leftmargin=2em]
  \item \textbf{Inputs \& Types:} \hrulefill\\[3mm]
  \item \textbf{Validation Criteria:} \hrulefill\\[3mm]
  \item \textbf{Outputs:} \hrulefill
\end{itemize}

\vspace{0.5mm}
\textbf{Task 2: Write the Structured Decision Logic in Ordered Steps} (7 points)\\
Formulate sequential pseudo-logic showing the order of validation, tiered branching, surcharges, and discount.

\begin{answerbox}[height=5.2cm]{Program-Logic Work Area}
\end{answerbox}

\vspace{0.5mm}
\textbf{Task 3: Hand Trace / Test Case Calculation} (6 points)\\
For $w = 6.0\text{ kg}$, \texttt{service = "E"}, and \texttt{is\_member = True}, show step-by-step arithmetic and final payable cost.

\begin{answerbox}[height=3.5cm]{Calculation Box for Task 3}
\end{answerbox}

\vspace{0.5mm}
\textbf{Task 4: Explain Two Design Decisions} (6 points)
\begin{enumerate}[label=\alph*),leftmargin=2em]
  \item Why is an \texttt{if-elif-else} ladder mandatory rather than consecutive independent \texttt{if} statements for tiered pricing?\\[3.5mm]
  \hrulefill
  \item Why must input validation be executed \textbf{prior} to calculating tiered prices or surcharges?\\[3.5mm]
  \hrulefill
\end{enumerate}

\newpage

% ==============================================================================
% PAGE 5: SECTION D, PROBLEM 1
% ==============================================================================

\section*{Section D: Applied Programming in Python \hfill \normalsize\normalfont(30 points)}
\textit{Write complete, syntactically correct Python statements. Apply proper conditional statements (\texttt{if}, \texttt{elif}, \texttt{else}), boolean operators, and input validation. Do \textbf{not} define custom functions (\texttt{def}) or loops. Suggested time: 18 minutes.}

\subsection*{Problem 1: Residential Electricity Billing with Progressive Tiered Tariff (15 points)}
Write a Python program that calculates the monthly electricity bill under a progressive tiered tariff structure with an excess surcharge and VAT.

\textbf{Requirements:}
\begin{enumerate}[leftmargin=2em,label=\arabic*.,itemsep=1pt]
  \item Prompt for the customer name \texttt{customer\_name} (string).
  \item Prompt for previous meter reading \texttt{prev\_reading} and current reading \texttt{curr\_reading} in kWh (convert to \texttt{float}).
  \item \textbf{Input Validation:} Both readings must be $\ge 0$, and $\text{curr\_reading} \ge \text{prev\_reading}$. If invalid, display an error message and do not compute the bill.
  \item Compute electricity consumed: $\text{kwh} = \text{curr\_reading} - \text{prev\_reading}$.
  \item Calculate the progressive energy charge: Tier 1 (\$0.12/kWh for first 50 kWh), Tier 2 (\$0.15/kWh for next 50 kWh, 51--100), Tier 3 (\$0.20/kWh for next 100 kWh, 101--200), and Tier 4 (\$0.25/kWh for consumption exceeding 200 kWh).
  \item \textbf{Excess Surcharge:} If total consumption exceeds 200 kWh, add a flat surcharge of \$10.00.
  \item \textbf{VAT Calculation:} Compute 8\% VAT on the subtotal (energy charge + surcharge).
  \item Display formatted receipt: Customer Name, Consumed kWh (1 decimal), Energy Charge, Surcharge, VAT (8\%), and Total Bill (all monetary values formatted to 2 decimals with \texttt{\$}).
\end{enumerate}

The output should follow this format when valid:
\begin{lstlisting}
Customer Name: Nguyen Van A
Electricity Consumed: 240.0 kWh
Energy Charge: $43.50
Excess Surcharge: $10.00
VAT (8%): $4.28
Total Bill: $57.78
\end{lstlisting}

\begin{answerbox}[height=12.2cm]{Python Solution for Problem 1}
\end{answerbox}

\newpage

% ==============================================================================
% PAGE 6: SECTION D, PROBLEM 2
% ==============================================================================

\subsection*{Problem 2: University Course Enrollment \& Eligibility Verification (15 points)}
Write a Python program that validates student eligibility to register for an advanced course (\texttt{DSAI2001 - Machine Learning}).

\textbf{Requirements:}
\begin{enumerate}[leftmargin=2em,label=\arabic*.,itemsep=1pt]
  \item Prompt for: \texttt{student\_id} (string), \texttt{gpa} (float), \texttt{credits\_earned} (int), \texttt{prereq\_status} (\texttt{"pass"}/\texttt{"fail"}), and \texttt{warning\_flag} (\texttt{"yes"}/\texttt{"no"}).
  \item \textbf{Student ID Format Validation:} ID must be 8 characters, start with \texttt{"NEU"}, and last 5 characters must be digits (\texttt{student\_id[3:].isdigit()}). If invalid, display: \texttt{"ERROR: Invalid Student ID format."} and terminate.
  \item \textbf{Eligibility Evaluation:} Eligible if and only if \texttt{prereq\_status.lower() == "pass"}, \texttt{warning\_flag.lower() == "no"}, $\text{credits\_earned} \ge 30$, and $\text{gpa} \ge 2.50$.
  \item \textbf{Enrollment Classification:}
    If eligible and $\text{gpa} \ge 3.60$ and $\text{credits\_earned} \ge 60$: status is \texttt{"APPROVED (Honors Priority)"}.
    If eligible otherwise: status is \texttt{"APPROVED (Standard)"}.
    If ineligible: status is \texttt{"REJECTED"}, stating the reason (e.g. missing prerequisite, low GPA, or warning).
  \item Display the formatted result card bordered with divider lines (\texttt{"=" * 45}).
\end{enumerate}

For example, given \texttt{NEU12345}, \texttt{3.75}, \texttt{64}, \texttt{pass}, and \texttt{no}:
\begin{lstlisting}
=============================================
ENROLLMENT ELIGIBILITY REPORT: DSAI2001
---------------------------------------------
Student ID: NEU12345 | Cumulative GPA: 3.75
Completed Credits: 64
Enrollment Status: APPROVED (Honors Priority)
=============================================
\end{lstlisting}

\begin{answerbox}[height=12.0cm]{Python Solution for Problem 2}
\end{answerbox}

\vfill
\begin{center}
  \rule{0.5\textwidth}{0.4pt}\\[1mm]
  \textbf{\color{primaryblue}--- END OF EXAM PAPER ---}
\end{center}

\end{document}
